%-------------------------------------------------------------

\medskip
\begin{enumerate}
	\item $f(-4)=(-4-1)(-4+3)=5$.
	\item On cherche à résoudre l'équation :
	\begin{align*}
	g(x)&=2\\
	2x+1&=2\\
	2x&=2-1\\
	x&=\dfrac{1}{2}
	\end{align*}
	L'antécédent de 2 par la fonction $g$ est $\dfrac{1}{2}$.
	\item \begin{enumerate}
				\item On doit saisir \og =2*B1+1 \fg{}.
				\item Pour $x=2$, on a $f(x)=g(x)$.
			\end{enumerate}
	\item \begin{enumerate}
				\item la fonction $g$ est une fonction affine donc sa représentation graphique est une droite, ainsi $\mathcal{C}_2$ est la courbe représentative de la fonction $g$ et $\mathcal{C}_1$ est la courbe représentative de la fonction $f$.
				\item Par lecture graphique (on regarde les abscisses des points d'intersection des deux courbes), on trouve $x=-2$ et $x=2$.
			\end{enumerate}
	\item Résolvons l'équation :
	\begin{align*}
	f(x)&=g(x)\\
	(x-1)(x+3)&=2x+1 \text{ \textit{On développe le membre de gauche.}}\\
	x^2+2x-3&=2x+1\\
	x^2+2x-3-2x-1&=0
	x^2-4&=0 
	\end{align*}
	Donc l'équation $f(x)=g(x)$ à les mêmes solutions que l'équation $x^2-4=0$.\\
	
	En factorisant $x^2-4$, on obtient $(x+2)(x-2)$	
	Résolvons $(x+2)(x-2)=0$
	Or, si un produit est nul alors l'un de ses facteurs est nul, donc :
	Soit $x+2=0$ ou $x-2=0$, donc
	$x=-2$ ou $x=2$.\\
	Ainsi on retrouve les solutions de la questions 4)b).
\end{enumerate}


\bigskip

%-------------------------------------------------------------
